%! Sine and Cosine Function Graphs — Unit 6, Lesson 6.5 — Guided Notes
\documentclass[10pt]{article}
\usepackage{precalculus-article}
\usepackage{precalculus-boxes}
\newcommand{\TallMath}[1]{$\displaystyle #1\rule[-1.4em]{0pt}{3.2em}$}

\begin{document}

\pageheader{Unit 6, Lesson 6.5}{Guided Notes}
\namedateperiod

\vspace{0.1in}

\begin{objectivebox}
By the end of this lesson, I will be able to:
\begin{itemize}[leftmargin=1.5em, itemsep=4pt]
    \item Explain how the unit circle produces the graph of $f(\theta)=\sin\theta$ and
          the graph of $f(\theta)=\cos\theta$. \writeline
    \item Identify the domain, range, period, midline, maximum, minimum, and zeros of
          each function from its graph. \writeline
\end{itemize}
\end{objectivebox}

\vspace{0.1in}

\begin{vocabbox}
\termblanklong{Oscillation}
\termblanklong{Amplitude (informal)}
\termblanklong{Period}
\termblanklong{Midline}
\termblanklong{Maximum point}
\termblanklong{Minimum point}
\termblanklong{Zero of a function}
\end{vocabbox}

\vspace{0.1in}

\begin{hookbox}
\textbf{Question:} We know the $y$-coordinate of a point on the unit circle at each angle.
What does the graph of that $y$-coordinate \emph{versus the angle} look like?

\vspace{0.06in}
\writelines{2}
\end{hookbox}

\vspace{0.1in}

\begin{notesbox}{Building the Graph of $f(\theta)=\sin\theta$}
The output of $f(\theta)=\sin\theta$ is the \blank{3cm} of the unit-circle point at
angle $\theta$.

\vspace{6pt}
Complete the table, then describe the resulting curve in words.

\vspace{6pt}
\renewcommand{\arraystretch}{1.55}
\begin{tabular}{@{}|c|c|@{}}
    \hline
    $\theta$ & $\sin\theta$ \\\hline
    $0$ & \blank{2cm} \\
    $\dfrac{\pi}{6}$ & $\dfrac{1}{2}$ \\
    $\dfrac{\pi}{4}$ & \blank{2cm} \\
    $\dfrac{\pi}{3}$ & $\dfrac{\sqrt{3}}{2}$ \\
    $\dfrac{\pi}{2}$ & \blank{2cm} \\
    $\pi$ & $0$ \\
    $\dfrac{3\pi}{2}$ & \blank{2cm} \\
    $2\pi$ & \blank{2cm} \\\hline
\end{tabular}
\quad
\begin{minipage}[t]{0.52\linewidth}\vspace{0pt}
\textbf{Plot these ordered pairs} on the $\theta$-$y$ axes below, then connect them
smoothly.

\vspace{0.5in}
\textit{(coordinate axes drawn by teacher on board)}

\vspace{0.3in}
The curve is called a \blank{4cm}. As $\theta$ increases from $0$ to $2\pi$, the output
\blank{3cm} between $-1$ and $1$.
\end{minipage}
\end{notesbox}

\vspace{0.1in}

\begin{notesbox}{Building the Graph of $f(\theta)=\cos\theta$}
The output of $f(\theta)=\cos\theta$ is the \blank{3cm} of the unit-circle point at
angle $\theta$.

\vspace{6pt}
\renewcommand{\arraystretch}{1.55}
\begin{tabular}{@{}|c|c|@{}}
    \hline
    $\theta$ & $\cos\theta$ \\\hline
    $0$ & \blank{2cm} \\
    $\dfrac{\pi}{2}$ & \blank{2cm} \\
    $\pi$ & \blank{2cm} \\
    $\dfrac{3\pi}{2}$ & \blank{2cm} \\
    $2\pi$ & \blank{2cm} \\\hline
\end{tabular}
\quad
\begin{minipage}[t]{0.55\linewidth}\vspace{0pt}
How does the cosine graph compare to the sine graph in overall shape?
\writelines{2}

How does the starting value at $\theta=0$ differ?
\writeline
\end{minipage}
\end{notesbox}

\vspace{0.1in}

\begin{notesbox}{Key Features of $\sin\theta$ and $\cos\theta$}
\renewcommand{\arraystretch}{1.7}
\begin{tabularx}{\linewidth}{@{} l X X @{}}
    \textbf{Feature} & \textbf{$f(\theta)=\sin\theta$} & \textbf{$f(\theta)=\cos\theta$} \\\hline
    Domain   & \blank{3.5cm} & \blank{3.5cm} \\
    Range    & \blank{3.5cm} & \blank{3.5cm} \\
    Period   & \blank{3.5cm} & \blank{3.5cm} \\
    Midline  & \blank{3.5cm} & \blank{3.5cm} \\
    Maximum value & \blank{3.5cm} & \blank{3.5cm} \\
    Minimum value & \blank{3.5cm} & \blank{3.5cm} \\
    Zeros on $[0,2\pi)$ & \blank{3.5cm} & \blank{3.5cm} \\
\end{tabularx}
\end{notesbox}

\vspace{0.1in}

\begin{practicebox}
Use the description of the sine graph over $[0,2\pi]$ to answer the following.

\begin{enumerate}[leftmargin=1.5em, itemsep=10pt]
    \item At what $\theta$-values in $[0,2\pi]$ does $\sin\theta=0$?

          \blank{8cm}

    \item At what $\theta$-value in $[0,2\pi]$ does $\sin\theta=1$?

          \blank{6cm}

    \item On what interval(s) in $[0,2\pi]$ is $\sin\theta$ increasing?

          \blank{8cm}

    \item On what interval(s) in $[0,2\pi]$ is $\sin\theta$ decreasing?

          \blank{8cm}
\end{enumerate}
\end{practicebox}

\end{document}
